%! Linear Functions: Slope, Forms & Graphing — Unit 1, Lesson 1.1 — Homework
% Gradual release: the "you do alone" block. Homework is generated per lesson (user direction,
% 2026-08-31); the teacher may substitute a DeltaMath set where the content is available there.
\documentclass[10pt]{article}
\usepackage{algebra2-article}
\usepackage{algebra2-boxes}
\usepackage{grading}
\gradingassignment{alg2-u01-l01-hw}
\gradingcourse{algebra_2}
\usepackage{algebra2-graded}
% Coordinate grid: {xmin}{xmax}{ymin}{ymax}
\newcommand{\lgrid}[4]{%
  \draw[step=1, linegray!40, very thin] (#1,#3) grid (#2,#4);
  \draw[->, semithick, charcoal] (#1,0)--(#2,0) node[below right, font=\scriptsize]{$x$};
  \draw[->, semithick, charcoal] (0,#3)--(0,#4) node[left, font=\scriptsize]{$y$};}

\begin{document}

\pageheader{Unit 1, Lesson 1.1}{Homework}

\vspace{0.10in}

\begin{notesbox}{Practice}
Show your work. Remember: \emph{slope-intercept} hands you a rate and a start, \emph{point-slope} is
built from a slope and \emph{any} point, and \emph{standard} is the tidy integer version --- all
three can describe the same line.

\begin{enumerate}[leftmargin=1.9em, itemsep=11pt]
  \item \textbf{Slope from two points.} \\[3pt]
        Through $(-2,7)$ and $(4,-5)$: $m = \inblank{\blank{1.4cm}}{\gradedblank[id=q1a, type=short, points=1, answer={-2}, check=numeric]{1.4cm-2pt}{10pt}}$ \quad
        Through $(-3,4)$ and $(6,4)$: $m = \inblank{\blank{1.2cm}}{\gradedblank[id=q1b, type=short, points=1, answer={0}, check=numeric]{1.2cm-2pt}{10pt}}$ \quad
        Through $(5,1)$ and $(5,-3)$: $m = $ \inblank{\blank{2.0cm}}{\gradedblank[id=q1c, type=short, points=1, answer={undefined}, check=text]{2.0cm-2pt}{10pt}}

  \boxguard[18]
  \item \textbf{One line, two ways.} The line shown passes through $(0,-1)$ and $(3,5)$.
    \begin{center}
    \begin{minipage}[c]{0.32\linewidth}\centering
    \begin{tikzpicture}[scale=0.30]
      \lgrid{-3}{5}{-6}{7}
      \draw[very thick, forest] (-2.5,-6)--(4,7);
      \fill[charcoal] (0,-1) circle (4pt);
      \fill[charcoal] (3,5) circle (4pt);
    \end{tikzpicture}
    \end{minipage}\hfill
    \begin{minipage}[c]{0.64\linewidth}
    \begin{enumerate}[label=(\alph*), leftmargin=1.9em, itemsep=4pt]
      \item Slope: \inblank{\blank{1.2cm}}{\gradedblank[id=q2a, type=short, points=1, answer={2}, check=numeric]{1.2cm-2pt}{10pt}} \; Slope-intercept form: \inblank{\blank{2.8cm}}{\gradedblank[id=q2b, type=short, points=1, answer={y=2x-1}, check=expression]{2.8cm-2pt}{10pt}}
      \item Point-slope form, using $(3,5)$: \inblank{\blank{3.6cm}}{\gradedblank[id=q2c, type=short, points=1, answer={y-5=2(x-3)}, check=expression]{3.6cm-2pt}{10pt}}
      \item Are (a) and (b) the same line? \inblank{\blank{1.4cm}}{\gradedblank[id=q2d, type=short, points=1, answer={yes}, check=text]{1.4cm-2pt}{10pt}} \; Show it:
    \end{enumerate}
    \end{minipage}
    \end{center}
    \begin{work}
      y - 5 &= 2(x - 3) \\
      y - 5 &= 2x - 6 \\
          y &= 2x - 1
    \end{work}

  \item \textbf{Change its outfit.} Rewrite $y + 2 = -\tfrac34(x - 4)$ in slope-intercept form,
        then in standard form.
    \begin{work}
      y + 2 &= -\tfrac34(x-4) \\
      y + 2 &= -\tfrac34 x + 3 \\
          y &= -\tfrac34 x + 1 \\
         4y &= -3x + 4 \\
    3x + 4y &= 4
    \end{work}
    Slope-intercept: \inblank{\blank{3.0cm}}{\gradedblank[id=q3a, type=short, points=2, answer={y=-(3/4)x+1}, check=expression, display={$y=-\tfrac34x+1$}]{3.0cm-2pt}{10pt}} \quad Standard: \inblank{\blank{2.6cm}}{\gradedblank[id=q3b, type=short, points=2, answer={3x+4y=4}, check=expression]{2.6cm-2pt}{10pt}}

  \boxguard[18]
  \item \textbf{The parent, and the two odd ones.} Start from $y = x$ and end at $y = -3x + 2$.
    \begin{enumerate}[label=(\alph*), leftmargin=1.9em, itemsep=4pt]
      \item Steeper or flatter? \inblank{\blank{1.4cm}}{\gradedblank[id=q4a, type=short, points=1, answer={steeper}, check=text]{1.4cm-2pt}{10pt}} \; Flipped? \inblank{\blank{1.0cm}}{\gradedblank[id=q4b, type=short, points=1, answer={yes}, check=text]{1.0cm-2pt}{10pt}} \;
            Shifted? \inblank{\blank{1.4cm}}{\gradedblank[id=q4c, type=short, points=1, answer={up 2}, check=text, display={up $2$}]{1.4cm-2pt}{10pt}}
      \item Slope of $y = 4$: \inblank{\blank{1.2cm}}{\gradedblank[id=q4d, type=short, points=1, answer={0}, check=numeric]{1.2cm-2pt}{10pt}} \quad Slope of $x = 4$: \inblank{\blank{2.0cm}}{\gradedblank[id=q4e, type=short, points=1, answer={undefined}, check=text]{2.0cm-2pt}{10pt}}
      \item Which of those two \emph{cannot} be written as $y = mx + b$, and why? \gradedlines{2}{\gradedblank[id=q4f, type=rubric, points=2, display={$x=4$. Its slope is undefined, so there is no number to put in for $m$; and every point on it has $x=4$, which $y=mx+b$ can never force.}]{\linewidth}{\gradedlinesheight}}
    \end{enumerate}

  \boxguard[20]
  \item \textbf{Model it.} A candle burns at a steady rate. After $1$ hour it is $17$~cm tall;
        after $4$ hours it is $11$~cm tall.
    \begin{enumerate}[label=(\alph*), leftmargin=1.9em, itemsep=4pt]
      \item Find the rate, then the rule.
        \begin{work}
          m &= \frac{11 - 17}{4 - 1} \\
            &= \frac{-6}{3} \\
            &= -2 \\
          h(t) &= -2t + 19
        \end{work}
      \item The $y$-intercept means \rubricbox{\gradedblank[id=q5a, type=rubric, points=2, display={the candle started at 19 cm}]{\linewidth}{10pt}}\par
            After $6$ hours:
            \inblank{\blank{1.4cm}}{\gradedblank[id=q5b, type=short, points=1, answer={7}, check=numeric, display={$7$ cm}]{1.4cm-2pt}{10pt}} \; Burns out at \inblank{\blank{1.8cm}}{\gradedblank[id=q5c, type=short, points=1, answer={9.5}, check=numeric, display={$t=9.5$ h}]{1.8cm-2pt}{10pt}}
    \end{enumerate}

  \item \textbf{Multiple choice (SOL-style).} Which line is \emph{perpendicular} to
        $y = \tfrac12 x - 4$ and passes through $(2,-1)$?
    \begin{enumerate}[label=(\Alph*), leftmargin=2.0em, itemsep=1pt]
      \item \letterbox{\gradedblank[id=q6a, type=choice, points=1, answer={B}]{0.4in}{10pt}}$y = \tfrac12 x - 2$
      \item $y = -2x + 3$
      \item $y = 2x - 5$
      \item $y = -\tfrac12 x$
    \end{enumerate}
    Say in one sentence how you ruled out one of the others. \gradedlines{2}{\gradedblank[id=q6b, type=rubric, points=2, display={\textbf{B} is correct: $m=-2$ is the opposite reciprocal of $\tfrac12$, and $-2(2)+3=-1$. \textbf{A} has the \emph{same} slope --- that is parallel, not perpendicular.}]{\linewidth}{\gradedlinesheight}}
\end{enumerate}
\end{notesbox}

\boxguard[18]
\begin{extensionbox}
\textbf{Build and reason.}
\begin{enumerate}[label=(\alph*), leftmargin=1.9em, itemsep=6pt]
  \item Write the equation of the line through $(-4,3)$ that is \emph{parallel} to $2x - y = 5$.
        \begin{work}
          2x - y &= 5 \\
               y &= 2x - 5 \\
           y - 3 &= 2(x + 4) \\
               y &= 2x + 11
        \end{work}
  \item In class you wrote a line from one of its two points; a classmate started from the other and
        got an equation that looked different. Explain why point-slope gives the \emph{same} line no
        matter which point on the line you start from. \writelines{2}
\end{enumerate}
\end{extensionbox}

\begin{spiralbox}
\textbf{Coming next --- Lesson 1.2: Absolute Value Equations \& Inequalities.} Every line you wrote
today crosses zero exactly once. Next comes a rule that reaches the same value \emph{twice}, because
it measures \emph{distance} --- so one equation will have two answers.
\end{spiralbox}

\end{document}
